\chapter{Introdução}

A gestão acadêmica semestral em instituições de ensino superior envolve um conjunto amplo de decisões interdependentes, entre as quais a atribuição de disciplinas a professores ocupa papel central. Trata-se de um problema recorrente, sensível ao contexto e de natureza combinatória: a cada período, mudanças no quadro docente, na oferta de disciplinas ou nas regras de carga horária podem alterar significativamente o espaço de soluções viáveis. Na literatura de pesquisa operacional, esse tipo de problema é estudado há décadas sob denominações como \textit{class--faculty assignment} e \textit{course timetabling} \cite{carter_laporte_1998,gunawan_ng_ong_2008}, o que reflete tanto a sua ubiquidade institucional quanto a sua complexidade computacional.

A abordagem padrão para modelar problemas de alocação docente é a Programação Linear Inteira (PLI), em que decisões binárias representam atribuições entre professores e disciplinas, e a função objetivo quantifica a qualidade da alocação por meio de pesos que combinam preferência e compatibilidade de área \cite{wolsey_2020}. Restrições lineares codificam requisitos operacionais, como carga horária mínima anual e máxima semestral por professor e cobertura obrigatória de todas as disciplinas. No contexto de atribuição de docentes com preferências, trabalhos como \citeonline{al_yakoob_sherali_2006} e \citeonline{domenech_lusa_2016} demonstram a eficácia dessa formulação para capturar exigências institucionais de forma sistemática. Modelos de PLI são resolvidos por métodos exatos como o Branch-and-Bound (B\&B), implementado em solvers modernos como o SCIP \cite{gamrath_2020_scip}, que utiliza relaxações lineares como limites para guiar a enumeração e podar regiões improdutivas da árvore de busca \cite{vanderbei_2014}.

Apesar da expressividade dos modelos exatos, instâncias de maior porte ou com restrições mais detalhadas podem dificultar a obtenção da prova de otimalidade dentro de limites de tempo práticos: o B\&B pode explorar um número exponencial de nós sem convergir, devolvendo soluções com gap significativo. No contexto específico do DPD (Problema da Distribuição de Disciplinas), \citeonline{silva2024abordagem_exata_dpd} investigou uma abordagem exata via PLI, enquanto \citeonline{jhonatan_grasp_dpd} desenvolveu uma heurística construtiva baseada em GRASP. O presente trabalho avança sobre essas contribuições ao propor uma matheurística de melhoria baseada em Large Neighborhood Search, combinando a qualidade de soluções iniciais obtidas por GRASP com uma etapa de refinamento via sub-MIP. A integração de matheurísticas baseadas em programação matemática --- em particular, a LNS no esquema \textit{fix-and-relax} como heurística primal do SCIP --- permanecia inexplorada para o DPD, configurando a lacuna metodológica que motiva o presente trabalho.

Esta monografia propõe, implementa e avalia uma matheurística baseada em LNS integrada ao solver SCIP para o DPD. O problema é modelado como um PLI binário e resolvido por B\&B, enquanto a LNS atua como mecanismo de intensificação: a partir de uma solução incumbente, uma fração controlada das atribuições é liberada para reotimização em um subproblema auxiliar (sub-MIP) com limite de tempo, e, caso uma solução de melhor qualidade seja encontrada, ela é incorporada ao solver principal como nova incumbente \cite{pisinger_ropke_2010}. Três parâmetros internos são analisados sistematicamente --- taxa de destruição, sentido de ordenação dos candidatos e tempo máximo do sub-MIP --- por meio de uma abordagem \textit{one-factor-at-a-time}. O desempenho da LNS é então comparado com três estratégias alternativas: o B\&B puro com relaxação linear, a heurística GRASP \cite{jhonatan_grasp_dpd} e uma variante híbrida GRASP|LNS, na qual a LNS é aplicada como etapa de melhoria a partir de soluções construídas pelo GRASP \cite{boschetti_letchford_maniezzo_2023,maniezzo_stutzle_voss_2021}.

A distribuição de disciplinas é realizada a cada semestre e está sujeita a prazos institucionais que demandam respostas rápidas, mesmo quando a prova de otimalidade não é exigida. Nesse cenário, a capacidade de gerar incumbentes de alta qualidade em estágios iniciais do B\&B tem impacto direto na eficiência da busca: soluções melhores permitem podas mais agressivas da árvore, acelerando a convergência \cite{neto_hoshino_pedrotti_2023}. Matheurísticas do tipo \textit{fix-and-relax} exploram precisamente essa propriedade ao resolver, repetidamente, versões restritas do problema original com estrutura preservada da incumbente corrente.

O restante deste trabalho está organizado da seguinte forma. No Capítulo~2 é apresentada a fundamentação teórica, cobrindo conceitos de otimização inteira, heurísticas e matheurísticas relevantes ao estudo. O Capítulo~3 define o DPD formalmente e descreve tanto o modelo matemático exato quanto a matheurística LNS proposta. O Capítulo~4 apresenta os resultados experimentais e a análise computacional, incluindo a calibração dos parâmetros internos da LNS e as comparações de desempenho entre as abordagens. Por fim, o Capítulo~5 conclui o trabalho e discute limitações e possibilidades de trabalhos futuros.

